
\subsubsection{Execution Feedback Reinforcement Learning}

PPOCoder \citep{ShojaeeEtAl2023PPOCoder} pioneered execution-based RL for code generation, training policy models with rewards from automated test case pass/fail signals. Their work demonstrated substantial correctness improvements on MBPP benchmarks, validating that non-differentiable execution feedback can integrate into policy gradient optimization. Process-Supervised RL \citep{YeEtAl2025ProcessSupervised} extends this with line-by-line verification using compiler feedback. These methods excel at functional correctness but optimize a single objective, ignoring code quality and maintainability.

Our Phase 1 mechanism builds on PPOCoder's execution feedback paradigm, using execution weight dominance during early training to establish correctness foundations. However, we treat execution feedback as the first stage of a multi-phase curriculum, not the sole training signal.

\subsubsection{Human Feedback Alignment for Code}

The RLHF paradigm, established for general language model alignment and adapted for code generation, trains models to optimize human preference scores from pairwise comparisons or quality ratings. This captures subjective quality dimensions that automated metrics cannot evaluate. However, human annotation is expensive, limiting scalability, and human-only training may sacrifice functional correctness.

ProSec \citep{XuEtAl2024ProSec} applies human feedback to security alignment. SEAlign \citep{ZhangEtAl2025SEAlign} extends human feedback to multi-step software engineering tasks. While these methods demonstrate value of human oversight, they apply human feedback uniformly across training.

Our approach uses human feedback strategically in Phase 3 with increasing weight (0.$400\rightarrow 0.636)$, reserving expensive oversight for edge case refinement after correctness and quality foundations are established.

\subsubsection{Multi-Criteria Reward Modeling}

Themis \citep{PaulEtAl2026Themis} represents state-of-the-art in multi-criteria code reward modeling, training reward models on 350,000+ preference pairs across multiple quality dimensions. However, Themis operates as an offline ranking system, not dynamically adjusting which criteria dominate during training.

Curriculum-RLAIF \citep{LiEtAl2025CurriculumRLAIF} introduces curriculum learning to alignment, but schedules task difficulty rather than feedback modality. Their AI feedback weight remains constant throughout training.

Our tri-modal framework extends multi-criteria approaches by integrating three distinct feedback sources online during RL training with dynamic weight scheduling: execution weight decays (0.$800\rightarrow 0.182)$, AI weight peaks mid-training (max 0.545 at 50\%), and human weight increases late (0.$100\rightarrow 0.636)$.

\subsubsection{Positioning Our Contribution}

The field has progressed from single-feedback methods toward multi-criteria integration and curriculum learning. Yet no prior work explores curriculum over feedback modality---dynamically adjusting which signal type receives emphasis across training phases. Our mechanism validation demonstrates that $execution\rightarrow AI\rightarrow human$ scheduling is implementable, produces predicted weight patterns, and achieves phase-specific objectives.
